\documentclass[11pt]{article}
% Shared preamble for all daily exercises
% Update this file to change settings (padding, parindent, etc.) for all exercises

\usepackage[margin=0.75in]{geometry}
\usepackage{amsmath}
\usepackage{amssymb}

% Custom settings
\setlength{\parindent}{0pt}
\setlength{\parskip}{0.2cm}

\title{Daily Exercise}
\author{Dhruman Gupta}
\date{21st April}

\begin{document}

% Common header for daily exercises
% This file can be included in other LaTeX documents

\maketitle

\noindent
\textbf{Course:} CS-3330, Theory of Computation

\vspace{0.2cm}

\noindent
\textbf{Question:}\noindent 
Define $\Sigma_0^P = \Pi_0^P = P$. For $k \geq 0$, let

\[
\Sigma_{k+1}^P = NP^{\Sigma_k^P}, \qquad
\Pi_{k+1}^P = coNP^{\Sigma_k^P}, \qquad
\Delta_{k+1}^P = P^{\Sigma_k^P}.
\]

Show that:
\begin{itemize}
    \item $NP^{NP} = \Sigma_2^P$
    \item $P^{NP} = \Delta_2^P$
    \item $NP^{NP^{NP}} = \Sigma_3^P$
\end{itemize}

\textbf{Answer:}

We know that $\Sigma_1^P = NP$.

First, we show that $NP^{NP} = \Sigma_2^P$. Let $L \in NP^{NP}$. Then some NP machine $M$ with NP oracle $A$ decides $L$.

Since $A \in NP$, there exists a polynomial-time predicate $S$ such that

\[
q \in A \iff \exists z \; S(q,z).
\]

Fix an input $x$ and an accepting branch of $M$. It asks only polynomially many oracle queries. For each yes-answer query $q_i$, we can guess a witness $z_i$ such that $S(q_i,z_i)$ holds. For each no-answer query $q_j$, we have

\[
\neg(\exists z\; S(q_j,z)) \iff \forall z\; \neg S(q_j,z).
\]

So the whole accepting computation has the form

\[
\exists y_1 \forall y_2 \; R(x,y_1,y_2)
\]

where $R$ is polynomial-time decidable. Here $y_1$ encodes the nondeterministic choices of $M$ and the witnesses for the yes-answer queries, and $y_2$ ranges over possible witnesses for the no-answer queries.

Hence $L \in \Sigma_2^P$, and so

\[
NP^{NP} \subseteq \Sigma_2^P.
\]

Conversely, let $L \in \Sigma_2^P$. Then

\[
x \in L \iff \exists y \forall z \; R(x,y,z)
\]

for some polynomial-time predicate $R$. Rewrite this as

\[
x \in L \iff \exists y \; \neg(\exists z \; \neg R(x,y,z)).
\]

An $NP^{NP}$ machine can guess $y$ and ask whether

\[
\exists z \; \neg R(x,y,z)
\]

holds. It accepts iff the oracle says no. Thus $L \in NP^{NP}$, and so

\[
\Sigma_2^P \subseteq NP^{NP}.
\]

Therefore,

\[
NP^{NP} = \Sigma_2^P.
\]

Next, by definition,

\[
\Delta_{k+1}^P = P^{\Sigma_k^P}.
\]

putting $k = 1$ gives

\[
\Delta_2^P = P^{\Sigma_1^P}.
\]

and since $\Sigma_1^P = NP$, this becomes

\[
\Delta_2^P = P^{NP}.
\]

Hence,

\[
P^{NP} = \Delta_2^P.
\]

Finally, we show that $NP^{NP^{NP}} = \Sigma_3^P$. Since

\[
NP^{NP} = \Sigma_2^P,
\]

we get

\[
NP^{NP^{NP}} = NP^{\Sigma_2^P}.
\]

Now let $L \in NP^{\Sigma_2^P}$. Then some NP machine $M$ with oracle $A \in \Sigma_2^P$ decides $L$.

Since $A \in \Sigma_2^P$, there exists a polynomial-time predicate $S$ such that

\[
q \in A \iff \exists y \forall z \; S(q,y,z).
\]

Fix an input $x$ and an accepting branch of $M$. It asks only polynomially many oracle queries. For a yes-answer query $q$, we have

\[
\exists y \forall z \; S(q,y,z),
\]

so its witness can be guessed existentially. For a no-answer query $q$, we have

\[
\neg(\exists y \forall z \; S(q,y,z)) \iff \forall y \exists z \; \neg S(q,y,z).
\]

So no-answers contribute one universal block followed by one existential block. Together with the outer NP guess, this gives a formula of the form

\[
\exists u \forall v \exists w \; R'(x,u,v,w)
\]

for some polynomial-time predicate $R'$. Hence $L \in \Sigma_3^P$, so

\[
NP^{\Sigma_2^P} \subseteq \Sigma_3^P.
\]

Conversely, let $L \in \Sigma_3^P$. Then

\[
x \in L \iff \exists y \forall z \exists w \; R(x,y,z,w)
\]

for some polynomial-time predicate $R$. Rewrite this as

\[
x \in L \iff \exists y \; \neg(\exists z \forall w \; \neg R(x,y,z,w)).
\]

The statement

\[
\exists z \forall w \; \neg R(x,y,z,w)
\]

is a $\Sigma_2^P$ predicate. So an $NP^{\Sigma_2^P}$ machine can guess $y$ and ask whether

\[
\exists z \forall w \; \neg R(x,y,z,w)
\]

holds, accepting iff the oracle answers no.

Hence $L \in NP^{\Sigma_2^P}$, so

\[
\Sigma_3^P \subseteq NP^{\Sigma_2^P}.
\]

Therefore,

\[
NP^{\Sigma_2^P} = \Sigma_3^P.
\]

Since $NP^{NP} = \Sigma_2^P$, we get

\[
NP^{NP^{NP}} = \Sigma_3^P.
\]

Hence proved.

\end{document}
